\section*{Exercises on \S\arabic{exercisegroup}}
\phantomsection
\addcontentsline{toc}{section}{Exercises on \protect\S\arabic{exercisegroup}}

\(\mathfrak{J}1)\) a) Let \(g\) be a continuous real function for \(x \geqslant 0\) . Show that if \(g\) satisfies the two identities

\[
\sum_ {k = 0} ^ {p - 1} g \left(\frac {x + k}{p}\right) = g (x)
\]

\[
\sum_ {h = 0} ^ {q - 1} g \left(\frac {x + h}{q}\right) = g (x)
\]

it also satisfies the identity

\[
\sum_ {j = 0} ^ {p q - 1} g \left(\frac {x + j}{p q}\right) = g (x).
\]

\begin{enumerate}
\def\labelenumi{\alph{enumi})}
\setcounter{enumi}{1}
\tightlist
\item
  Deduce from this that if a function \(g\) has a continuous derivative for \(x \geqslant 0\) and satisfies the identity
\end{enumerate}

\[
\sum_ {k = 0} ^ {p - 1} g \left(\frac {x + k}{p}\right) = g (x)
\]

then it is of the form \(a(x-\frac{1}{2})\) , where a is a constant (remark that g satisfies the analogous identity where p is replaced by \(p^{n}\) ; let n tend to \(+\infty\) and deduce that \(g'(x)=\int_{0}^{1}g'(t)dt\) ). c) Conclude from b) that the \(\Gamma\) function is the only function having a continuous derivative for x\textgreater0, which satisfies equation (1) of VII, p.~305, and the multiplication formula (12) of VII, p.~318, for a value of p.

\begin{enumerate}
\def\labelenumi{\arabic{enumi})}
\setcounter{enumi}{1}
\tightlist
\item
  For every integer \(k > 1\) put \(S_k = \sum_{n=1}^{\infty} n^{-k}\) . Show that, for \(-1 < x \leqslant 1\) one has
\end{enumerate}

\[
\log \Gamma (1 + x) = - \gamma x + \sum_ {k = 2} ^ {\infty} (- 1) ^ {k} \frac {\mathrm{S} _ {k}}{k} x ^ {k}
\]

the series on the right being uniformly convergent on every compact subset of \(]-1,1]\) , and absolutely convergent for \(|x|<1\) .

\begin{enumerate}
\def\labelenumi{\arabic{enumi})}
\setcounter{enumi}{2}
\tightlist
\item
  Let s be a fixed real number; show that when t tends to \(+\infty\) or to \(-\infty\) one has
\end{enumerate}

\[
| \Gamma (s + i t) | \sim \sqrt {2 \pi} | t | ^ {s - 1 / 2} e ^ {- (\pi / 2) | t |}
\]

and

\[
\frac {\Gamma^ {\prime} (s + i t)}{\Gamma (s + i t)} \sim \log | t |.
\]

\begin{enumerate}
\def\labelenumi{\arabic{enumi})}
\setcounter{enumi}{3}
\tightlist
\item
  Let t be a fixed real number \(\neq 0\) ; show that as s tends to \(+\infty\)
\end{enumerate}

\[
| \Gamma (- s + i t) | \sim \sqrt {\frac {\pi}{2}} \left(s ^ {s + 1 / 2} e ^ {- s} \sqrt {\sinh^ {2} \pi t + \sin^ {2} \pi s}\right) ^ {- 1}
\]

(use the complements' relation).

\begin{enumerate}
\def\labelenumi{\arabic{enumi})}
\setcounter{enumi}{4}
\tightlist
\item
  Let \(x_{n}\) be the root of the equation \(\Gamma'(x) = 0\) in the interval \(] - n, -n + 1[\) . Show that
\end{enumerate}

\[
x _ {n} = - n + \frac {1}{\log n} + O \left(\frac {1}{(\log n) ^ {2}}\right)
\]

(use the complements' relation and formula (5) of VII, p.~316). Deduce that

\[
\Gamma (x _ {n}) \sim \frac {(- 1) ^ {n}}{\sqrt {2 \pi}} n ^ {- n - 1 / 2} e ^ {n + 1} \log n.
\]

\begin{enumerate}
\def\labelenumi{\arabic{enumi})}
\setcounter{enumi}{5}
\tightlist
\item
  Let \(V_{n}\) be the Vandermonde determinant \(\mathrm{V}(1,2,\dots,n)\) (Alg., III, p.~532). Show that
\end{enumerate}

\[
\log V _ {n} = \frac {n ^ {2}}{2} \log n - \frac {3 n ^ {2}}{4} + \left(\frac {1}{2} \log 2 \pi - \frac {1}{4}\right) n - \frac {1}{12} \log n + k + O \left(\frac {1}{n}\right)
\]

where k is a constant.

